\documentclass[12pt, a4paper]{article}

% ── Packages ──────────────────────────────────────────────────────────────────
\usepackage[a4paper, top=2.5cm, bottom=2.5cm, left=2.5cm, right=2.5cm]{geometry}
\usepackage{amsmath, amssymb, amsthm}
\usepackage{mathtools}
\usepackage{booktabs}
\usepackage{tabularx}
\usepackage{array}
\usepackage[table]{xcolor}
\usepackage{hyperref}
\usepackage{enumitem}
\usepackage{titlesec}
\usepackage{parskip}
\usepackage{microtype}
\usepackage{fancyhdr}
\usepackage{mdframed}
\setlength{\headheight}{13.6pt}

% ── Colors ────────────────────────────────────────────────────────────────────
\definecolor{iztechblue}{RGB}{0, 56, 101}
\definecolor{lightgray}{RGB}{245, 245, 245}
\definecolor{contributiongreen}{RGB}{0, 100, 60}

% ── Section Formatting ────────────────────────────────────────────────────────
\titleformat{\section}
  {\large\bfseries\color{iztechblue}}
  {\thesection.}{0.6em}{}[\titlerule]

\titleformat{\subsection}
  {\normalsize\bfseries\color{iztechblue!80}}
  {\thesubsection.}{0.5em}{}

% ── Header / Footer ───────────────────────────────────────────────────────────
\pagestyle{fancy}
\fancyhf{}
\renewcommand{\headrulewidth}{0.4pt}
\fancyhead[L]{\small\color{gray} Master Thesis Research Proposal}
\fancyhead[R]{\small\color{gray} \textit{Gökay Gülsoy — IZTECH, 2026}}
\fancyfoot[C]{\small\thepage}

% ── Theorem-like environments ─────────────────────────────────────────────────
\newmdtheoremenv[
  backgroundcolor=lightgray,
  linecolor=iztechblue,
  linewidth=1.5pt,
  topline=false, rightline=false, bottomline=false,
  innerleftmargin=10pt, innerrightmargin=10pt,
  innertopmargin=8pt, innerbottommargin=8pt
]{contribution}{Contribution}

\hypersetup{
  colorlinks=true,
  linkcolor=iztechblue,
  urlcolor=iztechblue,
  citecolor=iztechblue
}

% ══════════════════════════════════════════════════════════════════════════════
\begin{document}
% ══════════════════════════════════════════════════════════════════════════════

% ── Title Block ───────────────────────────────────────────────────────────────
\begin{center}
  {\LARGE\bfseries\color{iztechblue}
    Optimizing Privacy-Preserving LLM Inference\\[4pt]
    via Polynomial Approximations within the CKKS Scheme}\\[1.2em]
  {\large Master of Science Thesis — Research Proposal}\\[0.4em]
  {\normalsize
    Gökay Gülsoy \quad|\quad
    Department of Computer Engineering, IZTECH\\[0.3em]
    \textbf{Advisor:} Assoc.\ Prof.\ Emrah İnan \quad|\quad Advisor Meeting — April 2026}
\end{center}

\vspace{0.5em}
\noindent\rule{\textwidth}{1.5pt}
\vspace{0.3em}

% ── 1. MOTIVATION ─────────────────────────────────────────────────────────────
\section{Motivation and Problem Statement}

The deployment of Large Language Models (LLMs) through cloud-based Machine Learning as a Service (MLaaS) platforms requires users to transmit sensitive prompts in plaintext to third-party servers, creating severe privacy risks in domains such as healthcare, legal, and financial services.

\textbf{Fully Homomorphic Encryption (FHE)} offers a cryptographic solution: computations are performed directly on \emph{ciphertext}, so the server never observes the raw query or response. The \textbf{CKKS scheme}~\cite{cheon2017} is the natural choice for neural network inference because it natively supports approximate arithmetic over real numbers — matching the numerical characteristics of floating-point deep learning.

\textbf{The core obstacle} is that modern Transformer architectures rely on three non-polynomial functions — \textbf{Softmax}, \textbf{GELU}, and \textbf{LayerNorm} — which cannot be evaluated natively in the encrypted domain. Every such function must be replaced by a polynomial approximation, and the degree of that polynomial directly determines the \emph{multiplicative depth} consumed from the finite CKKS level budget.

\textbf{The central tension of the thesis:}
\begin{itemize}[leftmargin=2em]
  \item Higher-degree polynomials $\;\Rightarrow\;$ better approximation accuracy, more depth consumed
  \item Lower-degree polynomials $\;\Rightarrow\;$ shallower circuits, lower memory/latency, worse accuracy
  \item Current literature uses uniform, fixed-degree approximations — the same polynomial at every layer
\end{itemize}

% ── 2. RELATED WORK AND GAP ───────────────────────────────────────────────────
\section{State of the Art and Identified Research Gaps}

\begin{table}[ht]
\centering
\caption{Summary of existing secure Transformer inference frameworks.}
\label{tab:related}
\renewcommand{\arraystretch}{1.3}
\begin{tabularx}{\textwidth}{lXXl}
\toprule
\textbf{Work} & \textbf{Non-linear Strategy} & \textbf{Key Limitation} & \textbf{Acc.\ Drop} \\
\midrule
CryptoNets~\cite{cryptonets} & $x^2$ activation (CNN only) & No Transformer support & $<1\%$ \\
Iron~\cite{iron}             & Fixed-interval minimax poly & Same poly at every layer; no data-aware intervals & 1--2\% \\
THE-X~\cite{thex}            & ReLU replacement + learned Softmax net & Client interaction required; no formal error bounds & 1.48\% \\
THEF~\cite{thef}             & HE + TEE hybrid (exact non-lin.\ in TEE) & Hardware trust assumption; side-channel risk & $\sim$0\% \\
AutoFHE~\cite{autofhe}       & Automated degree selection (CNNs) & CNNs only; not extended to Transformers & varies \\
BOLT~\cite{bolt}             & Optimized CKKS Transformer & State-of-the-art latency; fixed per-function degree & 1--3\% \\
\bottomrule
\end{tabularx}
\end{table}

\medskip
\noindent\textbf{Identified gaps:}\\[2pt]
\textbf{Gap 1 — Data-agnostic intervals.} All existing polynomial approximations use a fixed interval (e.g.\ $[-5, 5]$) regardless of the actual activation distribution of the model, wasting approximation capacity on tail regions that are rarely activated.\\[2pt]
\textbf{Gap 2 — Uniform depth allocation.} Every Transformer layer is assigned the same polynomial degree for each non-linear function, even though activation distributions—and thus the optimal degree-accuracy trade-off—vary significantly across layers.\\[2pt]
\textbf{Gap 3 — No formal optimality criterion.} THE-X's learned Softmax estimation network provides no formal error guarantee. AutoFHE addresses degree selection for CNNs but has not been extended to the more complex Transformer architecture.

% ── 3. PROPOSED CONTRIBUTIONS ────────────────────────────────────────────────
\section{Proposed Novel Contributions}

\begin{contribution}[Distribution-Weighted Minimax Polynomial Approximation]
Standard minimax minimizes worst-case error over a fixed geometric interval:
\begin{equation}
  p_n^* = \arg\min_{p \in \mathcal{P}_n}\ \max_{x \in [a,b]} \bigl|f(x) - p(x)\bigr|.
  \label{eq:standard_minimax}
\end{equation}
We instead solve a \textbf{density-weighted} variant:
\begin{equation}
  p_n^{**} = \arg\min_{p \in \mathcal{P}_n}\ \max_{x \in [a,b]}\ \rho(x)\cdot\bigl|f(x) - p(x)\bigr|,
  \label{eq:weighted_minimax}
\end{equation}
where $\rho(x) > 0$ is the empirical activation density at position $x$, obtained by profiling the plaintext model on representative task data. The effective support of $\rho$ is typically far narrower than the assumed interval $[a,b]$; for BERT-Tiny on SST-2, GELU inputs concentrate within $[-2.5, 2.5]$ even though the literature uses $[-5, 5]$. This formulation is \emph{equivalent} to a standard weighted Chebyshev approximation, computable via a modified Remez algorithm with a Jacobi-weighted inner product, and admits closed-form error bounds:
\begin{equation}
  \bigl\|f - p_n^{**}\bigr\|_{\rho,\infty} \;\leq\; \bigl\|f - p_n^{*}\bigr\|_\infty \cdot \|\rho\|_\infty,
\end{equation}
showing that the weighted solution is never worse than the unweighted minimax on the support of $\rho$.
\end{contribution}

\begin{contribution}[Layer-Adaptive CKKS Depth Allocation]
Given a global depth budget $D_{\mathrm{budget}}$ over $L$ Transformer layers, define the degree assignment vector $\mathbf{n} = (n_i^{(\ell)})$ where $i \in \{\mathrm{GELU},\,\mathrm{Softmax},\,\mathrm{LN}\}$ and $\ell = 1,\ldots,L$. The optimal allocation solves:
\begin{equation}
  \min_{\mathbf{n}\,\in\,\mathbb{Z}_{+}^{3L}}
  \sum_{\ell=1}^{L} \sum_{i}
  \mathcal{E}\!\left(n_i^{(\ell)},\, \rho_i^{(\ell)}\right)
  \quad\text{subject to}\quad
  \sum_{\ell=1}^{L} \sum_{i} \left\lceil \log_2 n_i^{(\ell)} \right\rceil \;\leq\; D_{\mathrm{budget}},
  \label{eq:depth_opt}
\end{equation}
where $\mathcal{E}(n, \rho)$ is the task-weighted approximation error for degree $n$ under empirical distribution $\rho_i^{(\ell)}$, and $\lceil \log_2 n \rceil$ is the multiplicative depth consumed by a degree-$n$ polynomial evaluated via binary tree. Because $n_i^{(\ell)}$ is integer-valued and small ($\leq 8$ in practice), this is a small \emph{integer program} solvable via dynamic programming in $O(L \cdot |I| \cdot D_{\mathrm{budget}})$ time, making it trivial to compute. The key insight is that \textbf{shallower layers and task-irrelevant activation ranges} can be handled with lower-degree polynomials, freeing depth budget for layers where accuracy is more sensitive.
\end{contribution}

\medskip
\noindent\textbf{Combined pipeline:}
\begin{center}
\textit{Profile activations $\to$ Fit weighted-minimax polynomials per layer $\to$ Solve depth allocation $\to$ Approximation-aware fine-tuning $\to$ CKKS encrypted inference}
\end{center}

% ── 4. COMPARISON FRAMEWORK ───────────────────────────────────────────────────
\section{Experimental Comparison Framework}

We evaluate four configurations on \textbf{BERT-Tiny} (2 layers, $d=128$) using GLUE benchmark tasks, to be extended to larger models if resources permit.

\begin{table}[ht]
\centering
\caption{Planned experimental configurations.}
\label{tab:configs}
\renewcommand{\arraystretch}{1.3}
\begin{tabular}{clll}
\toprule
\textbf{Config} & \textbf{Method} & \textbf{Degree Strategy} & \textbf{Interval Strategy} \\
\midrule
B1 & Taylor (baseline) & Uniform, fixed & Symmetric $[-5,5]$ \\
B2 & Uniform minimax (Iron-style) & Uniform, fixed & Symmetric $[-5,5]$ \\
B3 & Uniform LS regression & Uniform, fixed & Profile-based \\
\rowcolor{green!8}
P  & \textbf{Ours} (Contributions 1+2) & Layer-adaptive & Profile-based, density-weighted \\
\bottomrule
\end{tabular}
\end{table}

\medskip
\noindent\textbf{Metrics:}
\begin{itemize}[leftmargin=2em]
  \item \textbf{Accuracy:} GLUE task scores (SST-2 accuracy, MRPC F1, STS-B Pearson correlation) — plaintext vs.\ approximated vs.\ encrypted
  \item \textbf{Approximation error:} $L^\infty$ and task-weighted $L^2$ error per function per layer
  \item \textbf{Depth consumed:} Total multiplicative levels per inference pass
  \item \textbf{Latency \& memory:} Wall-clock time and RAM/VRAM for encrypted inference in TenSEAL
\end{itemize}

% ── 5. IMPLEMENTATION PLAN ────────────────────────────────────────────────────
\section{Implementation Roadmap}

\begin{table}[ht]
\centering
\caption{Phased implementation schedule aligned with thesis Gantt chart.}
\label{tab:roadmap}
\renewcommand{\arraystretch}{1.3}
\begin{tabularx}{\textwidth}{llX}
\toprule
\textbf{Phase} & \textbf{Period} & \textbf{Deliverables} \\
\midrule
Phase A & Apr--May 2026 &
  Profile BERT-Tiny activation distributions on SST-2/GLUE;
  implement Taylor, minimax, LS, and weighted-minimax fitters;
  solve depth allocation integer program;
  generate approximation error figures \\
Phase B & May--Jul 2026 &
  Approximation-aware fine-tuning of BERT-Tiny for all four configurations;
  plaintext accuracy benchmarks (GLUE);
  Chapter 4 results (plaintext section) \\
Phase C & Aug--Nov 2026 &
  TenSEAL/CKKS encrypted FFN (Contribution 2 first);
  TenSEAL encrypted self-attention;
  full encrypted Transformer block \\
Phase D & Nov 2026--Mar 2027 &
  CKKS parameter tuning;
  latency and memory profiling;
  comparison against Iron/BOLT baselines;
  Chapter 4 complete \\
Phase E & Mar--Jul 2027 &
  Cloud GPU experiments (Colab A100);
  final analysis and thesis writing \\
\bottomrule
\end{tabularx}
\end{table}

% ── 6. LITERATURE BASE ────────────────────────────────────────────────────────
\section{Literature Base}

\renewcommand{\arraystretch}{1.2}
\begin{tabular}{p{0.42\textwidth} p{0.52\textwidth}}
\toprule
\textbf{Paper} & \textbf{Role in Thesis} \\
\midrule
Cheon et al.\ (ASIACRYPT 2017)~\cite{cheon2017} & CKKS scheme — mathematical foundation \\
CryptoNets (ICML 2016)~\cite{cryptonets} & First encrypted neural network; $x^2$ baseline \\
Iron (NeurIPS 2022)~\cite{iron} & Fixed-degree minimax for Transformers; primary baseline \\
THE-X (ACL 2022)~\cite{thex} & Approximation-aware fine-tuning workflow \\
THEF (TrustCom 2024)~\cite{thef} & HE+TEE hybrid; contrasting architecture \\
Han \& Ki (EUROCRYPT 2020)~\cite{han2020} & Optimal minimax poly for CKKS; mathematical basis for Contribution 1 \\
AutoFHE (USENIX 2023)~\cite{autofhe} & Automated degree selection (CNNs); gap we extend to Transformers \\
BOLT (IEEE S\&P 2024)~\cite{bolt} & SOTA encrypted Transformer; performance upper bound \\
Cheon et al.\ (ASIACRYPT 2021)~\cite{cheon2021} & Composite polynomial evaluation; potential Contribution 3 extension \\
HE Standard (2018)~\cite{hestandard} & CKKS parameter selection and security levels \\
\bottomrule
\end{tabular}

% ── 7. EXPECTED CONTRIBUTIONS SUMMARY ────────────────────────────────────────
\section{Summary of Expected Contributions}

\begin{enumerate}[leftmargin=2em, label=\textbf{C\arabic*.}]
  \item \textbf{Distribution-weighted minimax polynomial approximation} for FHE-compatible non-linear functions: a mathematically rigorous generalization of standard minimax that exploits empirical activation densities, with provable error bounds.

  \item \textbf{Layer-adaptive CKKS depth allocation} via integer programming: the first systematic framework for assigning non-uniform polynomial degrees across Transformer layers under a global depth budget, applicable to any CKKS-based secure inference system.

  \item \textbf{Empirical evaluation} on BERT-Tiny/GLUE demonstrating accuracy, depth, and latency improvements over Iron-style and THE-X-style baselines — providing a complete comparison framework for future work.
\end{enumerate}

% ── 8. PLAIN-LANGUAGE SUMMARY ─────────────────────────────────────────────────
\section{Plain-Language Summary: What We Do, Why, and How}

\subsection*{The Problem in One Sentence}
When a user sends a query to a cloud-hosted language model (e.g.\ ChatGPT-style service), the server sees the raw text — an unacceptable privacy risk for sensitive domains like healthcare, law, and finance.

\subsection*{Our Goal}
Enable a language model to process a user's query \textbf{without ever decrypting it}. The user encrypts their text locally, the server runs the model entirely on encrypted data, and only the user can decrypt the answer. The server learns \emph{nothing} about the query or the response.

\subsection*{Why Is This Hard?}
Encryption schemes that allow computation on encrypted data (called \emph{Fully Homomorphic Encryption}, specifically the CKKS scheme) support only addition and multiplication — no divisions, exponentials, or square roots. Modern language models (Transformers) rely on three operations that use exactly these unsupported functions:

\begin{enumerate}[leftmargin=2em]
  \item \textbf{GELU} — the activation function inside each feed-forward block (uses exponentials)
  \item \textbf{Softmax} — computes attention weights (uses exponentials and division)
  \item \textbf{LayerNorm} — normalises each layer's output (uses division and square root)
\end{enumerate}

The only way to handle these under encryption is to \textbf{replace them with polynomials} (which are just additions and multiplications). But there is a catch: every multiplication under encryption costs a limited resource called \emph{depth}. Higher-degree polynomials give better accuracy but burn more depth. When depth runs out, the ciphertext breaks.

\subsection*{What Existing Work Does (and What It Misses)}

\renewcommand{\arraystretch}{1.4}
\noindent
\begin{tabular}{@{} p{0.38\textwidth} p{0.58\textwidth} @{}}
\toprule
\textbf{Current practice} & \textbf{What we improve} \\
\midrule
Same polynomial degree at every layer &
  We assign \textbf{different degrees to different layers} — layers that barely affect accuracy get cheaper polynomials, saving budget for the layers that matter most. \\[6pt]
Fixed approximation interval (e.g.\ $[-5,5]$) regardless of actual data &
  We \textbf{profile real activation values} from the model and focus the polynomial's accuracy on the range that actually occurs, instead of wasting capacity on values the model never produces. \\[6pt]
No principled way to split a limited depth budget &
  We formulate the depth-allocation as a \textbf{simple optimisation problem} (solvable in milliseconds) that automatically decides the best degree for each function at each layer. \\
\bottomrule
\end{tabular}

\subsection*{Our Two Contributions — In Plain Words}

\begin{mdframed}[backgroundcolor=green!5, linecolor=contributiongreen, linewidth=2pt,
  topline=false, rightline=false, bottomline=false,
  innerleftmargin=12pt, innerrightmargin=12pt,
  innertopmargin=10pt, innerbottommargin=10pt]
\textbf{Contribution 1 — Smarter Polynomials (Distribution-Weighted Minimax Approximation)}

\medskip
\textit{Analogy:} Imagine fitting a suit. Off-the-rack (existing work) gives everyone the same size. We take measurements first (profile activations), then tailor the suit (polynomial) to the wearer (actual data distribution). The result fits better where it matters, and we can \emph{prove} it is never worse than the off-the-rack version.

\medskip
\textbf{Technique:}
\begin{enumerate}[leftmargin=1.5em, itemsep=2pt]
  \item Run the plaintext model on task data and record every input to GELU, Softmax, and LayerNorm at each layer.
  \item Build a density profile: where do values concentrate? (e.g.\ GELU inputs mostly land in $[-2.5, 2.5]$, not $[-5, 5]$).
  \item Fit the polynomial using a \emph{weighted} minimax algorithm that penalises errors in high-density regions more than errors in tails.
\end{enumerate}

\medskip
\textbf{Status:} Fully implemented and tested. Our profiling of BERT-Tiny on SST-2 confirms that activation ranges are 30--50\% narrower than literature defaults.
\end{mdframed}

\vspace{0.8em}

\begin{mdframed}[backgroundcolor=green!5, linecolor=contributiongreen, linewidth=2pt,
  topline=false, rightline=false, bottomline=false,
  innerleftmargin=12pt, innerrightmargin=12pt,
  innertopmargin=10pt, innerbottommargin=10pt]
\textbf{Contribution 2 — Smarter Depth Budget (Layer-Adaptive Depth Allocation)}

\medskip
\textit{Analogy:} A university has a limited number of scholarship slots. Uniform allocation gives every department the same number. Our approach gives more slots to departments with the highest demand and impact, and fewer to departments that can manage with less.

\medskip
\textbf{Technique:}
\begin{enumerate}[leftmargin=1.5em, itemsep=2pt]
  \item For each function at each layer, pre-compute the approximation error at every candidate degree (2, 4, 6, 8, \ldots) using the density-weighted polynomial from Contribution 1.
  \item Define a total depth budget (determined by CKKS encryption parameters).
  \item Run a dynamic programming solver that distributes the budget across all functions/layers to minimise total error.
\end{enumerate}

\medskip
\textbf{Status:} Fully implemented. On BERT-Tiny with a budget of 15 depth levels, our adaptive allocation achieves \textbf{25.7$\times$ lower total error} than uniform allocation. The solver identifies Softmax at Layer 1 as the accuracy bottleneck (wide interval, high error), automatically assigning it the highest degree.
\end{mdframed}

\subsection*{Techniques Used and To Be Used}

\begin{center}
\renewcommand{\arraystretch}{1.35}
\begin{tabularx}{\textwidth}{>{\bfseries}l l X}
\toprule
\textbf{Technique} & \textbf{Status} & \textbf{Purpose} \\
\midrule
Activation profiling (hook-based) & \textcolor{contributiongreen}{Done} & Record real value ranges from BERT-Tiny on SST-2 \\
KDE density estimation & \textcolor{contributiongreen}{Done} & Build smooth density functions from recorded values \\
Chebyshev polynomial fitting & \textcolor{contributiongreen}{Done} & Compute optimal polynomial coefficients \\
Weighted Remez algorithm & \textcolor{contributiongreen}{Done} & Solve density-weighted minimax problem \\
Dynamic programming solver & \textcolor{contributiongreen}{Done} & Optimally allocate depth budget across layers \\
HuggingFace Trainer fine-tuning & \textcolor{orange!80!black}{In progress} & Baseline and poly-replaced BERT-Tiny on SST-2 \\
Model surgery (activation swap) & \textcolor{orange!80!black}{In progress} & Replace GELU/Softmax/LayerNorm with polynomial modules \\
Approximation-aware fine-tuning & \textcolor{gray}{Planned} & Re-train with polynomials inserted to recover accuracy \\
CKKS encrypted inference (TenSEAL) & \textcolor{gray}{Planned} & Run full inference on encrypted data \\
Latency/memory benchmarking & \textcolor{gray}{Planned} & Measure real-world cost of encrypted inference \\
\bottomrule
\end{tabularx}
\end{center}

\subsection*{What We Expect to Show}

\begin{itemize}[leftmargin=2em, itemsep=4pt]
  \item \textbf{Accuracy:} With approximation-aware fine-tuning, the polynomial-replaced model should recover to within \textasciitilde1\% of the original BERT-Tiny accuracy on SST-2 — \emph{competitive with or better than} Iron and THE-X baselines.
  \item \textbf{Depth savings:} The adaptive allocation should require 20--40\% fewer depth levels than uniform allocation for the same accuracy target, directly translating to smaller CKKS parameters, less memory, and lower latency.
  \item \textbf{Generalisability:} The framework (profile $\to$ fit $\to$ allocate $\to$ fine-tune) applies to any CKKS-based Transformer system, not just BERT-Tiny.
\end{itemize}

\vspace{0.5em}

% ── References ────────────────────────────────────────────────────────────────
\medskip
\noindent\rule{\textwidth}{0.5pt}
{\footnotesize
\begin{thebibliography}{99}
\bibitem{cheon2017} J.H.\ Cheon, A.\ Kim, M.\ Kim, Y.\ Song, ``Homomorphic Encryption for Arithmetic of Approximate Numbers,'' \textit{ASIACRYPT 2017}, LNCS 10624, pp.\ 409--437, Springer, 2017.
\bibitem{cryptonets} R.\ Gilad-Bachrach et al., ``CryptoNets: Applying Neural Networks to Encrypted Data with High Throughput and Accuracy,'' \textit{ICML 2016}, PMLR 48, pp.\ 201--210, 2016.
\bibitem{iron} M.\ Hao et al., ``Iron: Private Inference on Transformers,'' \textit{NeurIPS 2022}, vol.\ 35, pp.\ 15718--15731, 2022.
\bibitem{thex} T.\ Chen et al., ``THE-X: Privacy-Preserving Transformer Inference with Homomorphic Encryption,'' \textit{Findings of ACL 2022}, pp.\ 3510--3520, 2022.
\bibitem{thef} Z.\ Li et al., ``THEF: A Privacy-Preserving Framework for Transformer Inference Leveraging HE and TEE,'' \textit{TrustCom 2024}, pp.\ 1693--1700, IEEE, 2024.
\bibitem{han2020} S.\ Han and J.\ Ki, ``Efficient Bootstrapping for Approximate Homomorphic Encryption with Non-Sparse Keys,'' \textit{EUROCRYPT 2021}, LNCS 12696, pp.\ 587--617, Springer, 2021.
\bibitem{autofhe} W.\ Ao et al., ``AutoFHE: Automated Adaption of CNNs for Efficient Evaluation over FHE,'' \textit{USENIX Security 2023}, 2023.
\bibitem{bolt} Z.\ Pang et al., ``BOLT: Privacy-Preserving, Accurate and Efficient Inference for Transformers,'' \textit{IEEE S\&P 2024}, 2024.
\bibitem{cheon2021} J.H.\ Cheon et al., ``Efficient Homomorphic Comparison Methods with Optimal Complexity,'' \textit{ASIACRYPT 2021}, LNCS 13090, pp.\ 221--256, Springer, 2021.
\bibitem{hestandard} M.\ Albrecht et al., ``Homomorphic Encryption Standard,'' homomorphicencryption.org, 2018.
\end{thebibliography}
}

\end{document}
